% =========================================================
\section{Partie commande et partie opérative}

\subsection{Introduction}

Dans un système automatisé, la \textbf{partie opérative} (PO) exécute les actions physiques, tandis que la \textbf{partie commande} (PC) décide de ces actions à partir des informations reçues.

Cette séparation est fondamentale : elle permet de distinguer clairement ce qui \textbf{décide} de ce qui \textbf{agit}.

\subsection{La partie opérative (PO)}

\subsubsection{Définition}

La partie opérative est l'ensemble des composants qui réalisent physiquement le travail sur la matière d'œuvre, sous l'ordre de la partie commande.

Elle agit sur la matière, la déplace, la transforme, la mesure ou la contrôle.

\subsubsection{Composants de la partie opérative}

On distingue généralement :

\begin{itemize}
    \item les \textbf{actionneurs} : moteurs, vérins, électroaimants ;
    \item les \textbf{préactionneurs} : contacteurs, relais, distributeurs, variateurs ;
    \item les \textbf{effecteurs} : outils, pinces, buses, tapis ;
    \item les \textbf{capteurs} : fins de course, détecteurs de proximité, capteurs de niveau ;
    \item les \textbf{organes de transmission} : arbres, courroies, engrenages, tuyaux.
\end{itemize}

\subsubsection{Rôle de la partie opérative}

La PO a pour fonction de :

\begin{itemize}
    \item recevoir les ordres de la PC ;
    \item convertir une énergie (électrique, pneumatique, hydraulique) en action mécanique ;
    \item agir sur la matière d'œuvre ;
    \item informer la PC de son état par l'intermédiaire des capteurs.
\end{itemize}

\paragraph{Exemple :} dans une porte automatique, la PO comprend le moteur, le réducteur, le bras mécanique et les capteurs de présence.

\subsection{La partie commande (PC)}

\subsubsection{Définition}

La partie commande est l'ensemble des composants qui traitent les informations et élaborent les ordres destinés à la partie opérative.

Elle constitue le « cerveau » du système automatisé.

\subsubsection{Composants de la partie commande}

Selon le type d'automatisme, la PC peut être réalisée par :

\begin{itemize}
    \item des \textbf{relais} et \textbf{contacteurs} (automatisme câblé) ;
    \item un \textbf{automate programmable industriel (API)} ;
    \item un \textbf{microcontrôleur} ou un calculateur industriel ;
    \item des \textbf{modules logiques} programmables.
\end{itemize}

\subsubsection{Rôle de la partie commande}

La PC a pour fonction de :

\begin{itemize}
    \item acquérir les informations d'entrée ;
    \item traiter ces informations selon une logique définie ;
    \item élaborer les ordres de sortie ;
    \item transmettre ces ordres à la PO ;
    \item gérer les sécurités, les alarmes et les modes de fonctionnement.
\end{itemize}

\paragraph{Exemple :} dans une porte automatique, la PC est un API qui reçoit l'information du capteur de présence et commande le moteur d'ouverture.

\subsection{Dialogue entre PC et PO}

La PC et la PO communiquent en permanence par deux flux d'informations :

\begin{itemize}
    \item la PC envoie des \textbf{ordres} à la PO ;
    \item la PO renvoie des \textbf{comptes rendus} à la PC.
\end{itemize}

On peut représenter ce dialogue par :

\[
\text{PC} \xrightarrow{\text{ordres}} \text{PO}
\qquad \text{et} \qquad
\text{PO} \xrightarrow{\text{comptes rendus}} \text{PC}
\]

\paragraph{Exemple :}

\begin{itemize}
    \item Ordre : « démarrer le moteur du convoyeur ».
    \item Compte rendu : « le convoyeur est en marche » ou « le carton est détecté ».
\end{itemize}

\subsection{Les informations échangées}

\subsubsection{Les ordres}

Les ordres sont élaborés par la PC. Ils peuvent être :

\begin{itemize}
    \item \textbf{TOR} (Tout Ou Rien) : marcher/arrêter, ouvrir/fermer ;
    \item \textbf{analogiques} : consigne de vitesse, de position, de température ;
    \item \textbf{numériques} : envoi d'un code, d'un compteur, d'une recette.
\end{itemize}

\subsubsection{Les comptes rendus}

Les comptes rendus sont fournis par les capteurs de la PO. Ils peuvent être :

\begin{itemize}
    \item \textbf{TOR} : présence, absence, fin de course ;
    \item \textbf{analogiques} : niveau, pression, température, débit ;
    \item \textbf{numériques} : position d'un codeur, comptage.
\end{itemize}

\subsection{Exemple complet : remplissage de cuve}

Considérons une cuve à remplir automatiquement.

\paragraph{Partie opérative :}
\begin{itemize}
    \item électrovanne de remplissage ;
    \item pompe d'alimentation ;
    \item capteurs de niveau bas et haut ;
    \item cuve.
\end{itemize}

\paragraph{Partie commande :}
\begin{itemize}
    \item API.
\end{itemize}

\paragraph{Dialogue :}

\[
\text{Niveau bas détecté} \rightarrow \text{API} \rightarrow \text{Ordre d'ouverture de l'électrovanne}
\]

\[
\text{Niveau haut détecté} \rightarrow \text{API} \rightarrow \text{Ordre de fermeture de l'électrovanne}
\]

\paragraph{Analyse :}
\begin{itemize}
    \item Le capteur de niveau bas est un \textbf{compte rendu} envoyé par la PO.
    \item L'API élabore un \textbf{ordre} d'ouverture ou de fermeture.
    \item L'électrovanne est un \textbf{actionneur} de la PO.
\end{itemize}

\subsection{Exemple complet : convoyeur}

\paragraph{Partie opérative :}
\begin{itemize}
    \item moteur du convoyeur ;
    \item tapis ;
    \item capteur de présence de carton ;
    \item contacteur.
\end{itemize}

\paragraph{Partie commande :}
\begin{itemize}
    \item API.
\end{itemize}

\paragraph{Dialogue :}

\[
\text{Carton détecté} \rightarrow \text{API} \rightarrow \text{Contacteur} \rightarrow \text{Moteur}
\]

\[
\text{Fin de convoyage} \rightarrow \text{API} \rightarrow \text{Arrêt du moteur}
\]

\subsection{Place de l'IHM}

L'IHM (Interface Homme-Machine) permet à l'opérateur de dialoguer avec la partie commande.

Elle peut :

\begin{itemize}
    \item envoyer des ordres à la PC (marche, arrêt, choix de mode) ;
    \item afficher les comptes rendus (états, alarmes, mesures) ;
    \item visualiser le fonctionnement du système.
\end{itemize}

Elle est reliée à la PC et non directement à la PO.

\subsection{Tableau récapitulatif}

\begin{table}[h]
\centering
\begin{tabularx}{\textwidth}{|X|X|}
\hline
\textbf{Partie commande (PC)} & \textbf{Partie opérative (PO)} \\
\hline
Traite les informations & Exécute les actions \\
\hline
API, relais, calculateur & Moteurs, vérins, contacteurs \\
\hline
Reçoit les comptes rendus & Fournit les comptes rendus \\
\hline
Élabore les ordres & Reçoit et exécute les ordres \\
\hline
Décide & Agit \\
\hline
\end{tabularx}
\caption{Répartition des rôles entre PC et PO}
\end{table}

\subsection{Conclusion}

La distinction entre partie commande et partie opérative est essentielle pour comprendre et concevoir un système automatisé.

La PC \textbf{décide}, la PO \textbf{agit}. Elles communiquent par des ordres et des comptes rendus. L'IHM permet à l'opérateur d'intervenir sans agir directement sur la PO.

\subsubsection*{À retenir}

\begin{quote}
La partie commande élabore les ordres à partir des informations reçues. La partie opérative exécute ces ordres sur la matière d'œuvre et renvoie des comptes rendus. L'IHM permet à l'opérateur de dialoguer avec la partie commande.
\end{quote}

\paragraph{Schéma fondamental :}

\[
\text{IHM} \rightarrow \text{Partie Commande} \rightarrow \text{Partie Opérative} \rightarrow \text{Matière d'œuvre}
\]

\[
\text{Ordres} \ \downarrow \qquad \uparrow \ \text{Comptes rendus}
\]